\section{Introduction}

The ability to generate multiple variants of an object during a construction process is  becoming increasingly frequent in many application areas. Most of the time, tools and operations used to create those variants are dedicated to specific fields and the construction process is often both tedious and time-consuming.

For example, in the field of Archaeology, remaining data found on the working field often represent only vestiges of ancients buildings. By means of 3D
reproduction, archaeologists painstakingly develop a number of hypotheses
they expect to test while being able to quickly model and
visualize them \cite{Quattrini_Baleani_2015}. CAD uses parametric history-based systems; such systems can be thought as dual structures with,  on the one hand, the geometric model corresponding to the modelled
object and,  on the other hand, the successive operations (and their parameters) recorded during the construction process. This process can then be reevaluated after some slight modifications upon the operations, without starting all over again from the beginning. Nevertheless, creating complex objects always requires a substantial amount of time.
Modeling buildings is also of interests for architects. Grammar-based procedural methods are commonly used
to generate variations of constructions \cite{haegler2009procedural, muller2006procedural}. But those grammars are based on a specific corpus of information which is difficult to transpose to other case studies.

In this paper, we propose to use  a rule-based graph transformation formalism, and more specifically the Jerboa software \cite{belhaouari2014jerboa}, to make the development of dedicated modelers for
specific applications easier. Rule-based languages for modeling
are commonly used in a number of fields such as plant growth with L-systems \cite{lindenmayer1974adding, bohl2015modeling}, wood’s internal structure \cite{terraz20093gmap}, or virtual cities 3D models
\cite{esri2023cityengine, beirao2010implementing}. Contrary to other approaches, Jerboa is independent of any specific
application field and does not require \textit{ad hoc} operations to
be manually coded. Operations are formally defined as rules within the Jerboa interface, thereby facilitating their rapid development. Furthermore, it guarantees the topological %coherence 
consistency of the underlying geometric model,
regardless of the applied operations \cite{arnould2022preserving}.

Jerboa is based on the generalized maps
(or G-maps) topological model \cite{lienhardt1991topological, Damiand-Lienhardt14}. This model
represents a specific class of labelled graph and allows the
homogeneous modeling of quasi-manifolds in any dimension. Number of
applications already make use of Jerboa and/or G-maps in
fields such as plant growth \cite{bohl2015modeling}, architecture \cite{horna2009consistency, arroyo2015dimension}, geology \cite{halbwachs2000generalized}  or
physics-based modeling  \cite{salah2017general}.

Despite its advantages, Jerboa does not provide the mechanisms required to quickly reevaluate variations from a base model. Conversely,  history-based parametric systems take advantage of the construction process recording to make reevaluations as fast and accurate as possible. Thus, our objective is
to extend the capabilities of Jerboa by incorporating the mechanisms inherent in parametric systems.

Any reevaluation of the parametric specification entails modifying the parameters of the operations. Those parameters are either geometric (such as the length of a groove) or references %to 
topological entities (vertices, edges, faces and so on) defined at an earlier stage of the modeling process.

Modifying some operation parameters requires ensuring that the subsequent operations are still valid, even if their own parameters have been updated. This issue, known as \textit{persistent naming}, is illustrated in Fig.~1.

\begin{figure}[t]
	\centering
	\begin{subfigure}[t]{\linewidth}
		\centering
		\includegraphics[width=.8\textwidth]{persistent-naming-pb-1}
		\caption{Initial evaluation}
		\label{fig:persistent-naming-pb-eval}
	\end{subfigure}

	\begin{subfigure}[t]{\linewidth}
		\centering
		\includegraphics[width=.8\textwidth]{persistent-naming-pb-2}
		\caption{Reevaluation}
		\label{fig:persistent-naming-pb-reval}
	\end{subfigure}
	\caption{Parametric specification}
	\label{fig:persistent-naming-pb}
\end{figure}

The initial parametric specification consists of three constructive operations (Fig.~\ref{fig:persistent-naming-pb-eval}):
\texttt{1-CreateCylinder(geo\_param1)}; \texttt{2-RoundedGroove(f1,geo\_param2)}; \texttt{3-Cy\-lin\-dri\-cal\-Pro\-tru\-sion(f2, geo\_param3)}.
When identifiers
or pointers (\textit{i.e.} concrete names) to the topological parameters of a geometric
model (\textit{e.g.} the identifier of the face \verb+f2+ as a parameter of
the \verb+CylindricalProtrusion+ operation) are used as topological
parameters, the issue of the persistence of these references at
reevaluation comes up. For example, in Fig.~\ref{fig:persistent-naming-pb-reval},
the rounded groove’s length is reduced. The
face \verb+f1+ is not split  anymore, unlike  during
the initial evaluation. Thus, neither face \verb+f2+ nor face \verb+f3+ are created:  identifiers and pointers to the entities are obviously different and, therefore, the cylindrical protrusion can
no longer be re-applied onto \verb+f2+.
Hence the necessity to use persistent
identifiers as operations parameters, which make possible to
unambiguously characterize entities and find their match at
reevaluation. In Fig.~\ref{fig:persistent-naming-pb}, using a persistent name to characterize
\verb+f2+ %from
during the initial evaluation  allows %to match
matching
it with the face \verb+fx+ at reevaluation.

Although persistent naming has been studied for decades in the CAD's
field \cite{KRIPAC1997,CAPOYLEAS199617,wu2001face,wang2005geometry,bidarra2005feature,mun05,mar06,baba10,xue16,farjana2018mechanisms},

% To our knowledge \highlight{only one preliminary approach} has attempted to use graph transformation rules to tackle this issue \cite{cardot2019persistent}.
\highlight{Two preliminary approaches have attempted to use graph transformation rules to tackle this issue \cite{cardot2019persistent,gaide2023model}.}
The authors' proposed approach in \cite{cardot2019persistent} to use \emph{History Records} (HRs) to represent the history of any topological entity designated in parametric specification and \emph{Matching Trees}(MT) to match this entity during reevaluation.
This is an interesting initial theoretical approach based on graph transformation, but the history represented in HR is limited.
In particular, no distinction is made between the entities at the origin of the designated entity and the evolution of the designated entity itself.
Furthermore, some elements are omitted in HRs (such as the history of entities at the origin of the designated entity), which can lead to mismatches during reevaluation.
\maxime{ajouter description de \cite{gaide2023model} avant de dire que l'on s'appuie dessus pour nos travaux.}
In this paper, we base our work on \cite{gaide2023model} to integrate the complete history of topological entities.
Our contribution is twofold. First, we widen the naming problem studies to the rule-based graph transformation modeling systems.
Second, we integrate the mechanisms of reevaluation for parametric systems into Jerboa.

We propose a persistent
naming method taking advantage of the %rules’ 
rule-based
formalization
of operations and their ability to precisely describe the history
of topological entities, such that these entities are uniquely and
homogeneously characterized for all dimensions.  Most existing methods require  tracking
numerous topological entities and consider the persistent naming problem only
through the prism of parameters modifications
from
a parametric specification standpoint \cite{KRIPAC1997, chen1995editability, wu2001face}. Our solution tracks
only the entities used in the parametric specification and the
ones they originate from. Moreover, not only the naming problem is tackled within the usual framework of  parameter edition, but we also take the specification edition (\textit{i.e.} adding, deleting and changing the order of operations) into account.

In section~2, we present the
main concepts necessary to carry out the persistent naming
mechanisms within the framework of a system, making use of
graph transformation rules. We
focus on G-maps, Jerboa’s rules, and how to
automatically detect topological
changes (creation, deletion, split, merging, modification) that
may occur upon any rule application. Section~3 is dedicated to the data structures
used by parametric specifications and persistent naming. Section~4 describes how
a parametric specification is evaluated or reevaluated through directed acyclic graphs which track the evolutions of topological entities and the %entities
ones
they originate from. Section~5 presents the matching process between evaluated and reevaluated entities. We conclude in section~6 and present the main directions of our future works.


%%% Local Variables:
%%% mode: latex
%%% TeX-master: "../main"
%%% End:
